\section{The S2 group: a joint three-piece estimate}
\label{sec:s2}

This section bounds the four S2 rows jointly from below. These rows are
$\{\text{S2-0},\text{S2-1},\text{S2-2},\text{S2-3(K/PF)}\}.$
Their signed lower functional is \texttt{S2-T8/4}.

\subsection{Source partition and four rows}
\label{sec:s2-1}\label{sec:s2-2}

\begin{lemma}
The plain S2 source satisfies (3.2), with the disjoint partition (3.4).
The selected third row has level and cutoff (3.5).
\end{lemma}
\begin{proof}
Fix
\[
 u=\frac1{500},\quad \beta=\frac{50}{309},\quad
 \epsilon_A=\eta_L=\frac1{10000},\quad
 0<\sigma\le10^{-10},\quad \tau=\sigma^9,\quad \rho=1+\tau,
\]
where $\sigma$ is sufficiently small for the AA37 constant and is fixed before $N$.
Then $N$ tends to infinity through even integers. Put
\[
 \mathcal A_D(v)=\mathbf1_{\{v\ {\rm prime}\}}
 \mathbf1_{\{N/\log^{30}N<v<(1-\epsilon_A)N\}},
 \qquad \mathcal A_0=\{N-v:\mathcal A_D(v)=1\}.
\]
Sections~3--5 use this endpoint, with the normalization
\[
 C(N)=\prod_{p>2}\left(1-\frac1{(p-1)^2}\right)
       \prod_{\substack{p\mid N\\p>2}}\frac{p-1}{p-2},\quad
 U_C=C(N)\frac{N}{\log^2N},\quad W_N=4(1-\epsilon_A)U_C.
 \tag{3.1}\label{p2:eq:s2-normalization}
\]
The first product is positive, and the second is at least one.
Thus $W_N\gg N/\log^2N$ uniformly in even $N$.

The unsifted AA37 input is $\mathcal A^{\rm pl}_{2,1}=\mathcal A_0$.
By (1.7),
$S_{2,\rm raw}=S(\mathcal A_0,\mathcal P_N,N^\beta)+\Delta^{\rm src}_2$,
where $\Delta^{\rm src}_2$ contains low-tested-prime and exceptional $v\mid N$ cells.
Source-square restoration is empty at $b=1$.
For a Buchstab prefix $b$ over $\mathcal P_N$, put
\[
 S_{\rm pl}(b,z)=S(\mathcal A^{\rm pl}_{2,b},\mathcal P_N,z),
 \qquad \mathcal A^{\rm pl}_{2,b}=\{N-v:\mathcal A_D(v)=1,\ b\mid N-v\},
\]
where $N-v$ is sieved with the empty cofactor allowed. Then
\[
\begin{split}
 T_0&=S_{\rm pl}(1,N^u),\\
 T_1&=\sum_{N^u\le p_1<N^\beta}S_{\rm pl}(p_1,N^u),\\
 T_2&=\sum_{N^u\le p_2<p_1<N^\beta}S_{\rm pl}(p_1p_2,N^u),\\
 T_3&=\sum_{N^u\le p_3<p_2<p_1<N^\beta}S_{\rm pl}(p_1p_2p_3,p_3),\\
 S_{\rm pl}(1,N^\beta)&=T_0-T_1+T_2-T_3.
\end{split}\tag{3.2}
\]
This follows by applying $S_{\rm pl}(b,z_2)=S_{\rm pl}(b,z_1)-\sum_{z_1\le p<z_2}S_{\rm pl}(bp,p)$ twice,
with repeated factors of $N-v$ still allowed.
The final cutoff is $p_3$.
Modulus-one coefficients give $(f,F,f,F)$, and the cutoff factors $(500,-500,500,-1)$
come from $1/u$ and $p_3$.

Let $r_i=\log_Np_i$. Put
$\Delta=\{u\le r_3<r_2<r_1<\beta\},\quad
 g(r)=r_1+2r_2+4r_3,\quad H=\frac{249}{250}.$
Set $PF=\Delta\cap\{g\le H\}$ and $R1=\Delta\cap\{g>H\}$.
The cells carry the inherited global half-open labels
$t_i\le r_i<\rho t_i,\qquad
 \mathcal P(t_i)=\{p\nmid2N:N^{\max(t_i,u)}\le p<N^{\min(\rho t_i,\beta)}\}.$
Clipping at $u$ or $\beta$ preserves each label.
A whole cell is selected when
\[
 h(t)=t_1+2t_2+(4+2\tau)t_3\le K,
 \qquad K=H+2\tau\beta.
 \tag{3.3}
\]
Let $\mathcal F$ be the selected union in $\Delta$.
Its strict complement $\mathcal B$ has $h>K$.
For $g\le H$, $h\le g+2\tau t_3\le K$,
and on $\mathcal F$, $g<\rho(t_1+2t_2+4t_3)\le\rho K$. Hence
\[
 PF\subseteq\mathcal F,\quad\mathcal B\subseteq R1,\quad
 \Delta=PF\ \dot\cup\ (\mathcal F\cap R1)\ \dot\cup\ \mathcal B,
 \quad \mathcal F\cap R1\subseteq\{H<g<\rho K\}.
 \tag{3.4}
\]
Equality belongs to PF or the selected label.
All three pieces retain the negative sign of $T_3$.
The coupled cut $g=H$ is compared with one whole-$K$ product-block estimate.

On a selected prefix block, put
\[
 L_b=N^{\Theta_b-\eta_L-\sum t_i},\qquad
 z=N^u\quad(r=0,1,2),\qquad z=E_3=N^{t_3}\quad(r=3),
 \tag{3.5}
\]
where $\Theta_b$ is the actual analytic cell infimum.
Since $E_3\le p_3$,
$S_{\rm pl}(b,p_3)\le S_{\rm pl}(b,E_3)$, the required upper bound for the negative row.
Its sieve argument exceeds $1.0117$.
The whole-$K$ cells satisfy the admission envelope $h\le1-2\eta_L$.
Here $K<1-2\eta_L,\qquad (1-2\eta_L)-H=19/5000.$
The boundary estimate retains the sharper condition (3.3).
\end{proof}

\subsection{Coefficients and prepared hypotheses}
\label{sec:s2-3}

\begin{lemma}
The four signed rows satisfy H1--H7 of Theorem~\ref{thm:distribution-input}.
Their coefficients are (3.6)--(3.8), with both masks (3.11).
\end{lemma}
\begin{proof}
Fix a complete symmetric prefix/tail block $k$ with its AA37 signs, and put
$C_k(q)=\sum_{d_{AA},(p_j)}\lambda_{d_{AA}}
                   \mathbf1_{\{q=d_{AA}\prod_jp_j\}}.$
The smooth support is
$d_{AA}<L_b^\sigma$, $d_{AA}\mid P_N(L_b^{\sigma^2})$.
It is squarefree, with primes avoiding $2N$.
The empty tail is included once.
Before unequal-coordinate dissection, complete equal-bin prefixes are divided by $M_k=\prod_G|G|!$.
Packing puts $d_{AA}$ wholly in slot 1, and the labelled convolution of slot $s$ is divided by $|J_s|!$. Thus
\[
 \frac{C_k}{M_k}=\frac{C_{\rm slot}}{M_k}
                         (\gamma_1*\gamma_2*\gamma_3),\quad
 C_{\rm slot}=\prod_{s=1}^3|J_s|!,\quad |\gamma_s|\le1.
 \tag{3.6}
\]
The prefix and slot factorials have distinct roles.
The signed remainder and its external scalar are
\[
 R_j=\sum_q(\gamma_1*\gamma_2*\gamma_3)(q)\Delta_j(q;N),
 \qquad c_j=w_kc_{{HB},{\rm box}}\frac{C_{\rm slot}}{M_k}.
 \tag{3.7}
\]
Prime realizations and smooth divisors stay inside the convolution.
Absolute values are taken after the signed $q$-sum.

The local density is $\omega_b(p^a)=p/(p-1)$ for every $a\ge1$, with higher-power constant 3.
Put $U_k=C_k/M_k$.
Let $\Xi_{\rm union}$ be the disjoint forbidden tail/prefix correction and
$\Xi_{\rm allowed}$ the allowed-repeat correction.
Let $W_k$ be the retained squarefree coefficient.
Then $O_k=U_k-\Xi_{\rm union}=W_k+\Xi_{\rm allowed}.$
With $m_A=\sum_v\mathcal A_D(v)$, define
\[
 \Delta_0(q)=\sum_{v\equiv N(q)}\mathcal A_D(v)-\frac{m_A}{\varphi(q)},
 \quad
 \Delta_P(q)=\sum_{v\equiv N(q)}\mathcal A_D(v)
             -\frac{\sum_{(v,q)=1}\mathcal A_D(v)}{\varphi(q)}.
\]
Substitution before HB gives
\[
 R_{AA}=S_{\rm slot}+NC_{sf}
 -\sum_{k,q}w_k\Xi_{\rm union}(q)\Delta_P(q)
 +\sum_{k,q}w_k\Xi_{\rm allowed}(q)(\Delta_0-\Delta_P)(q).
 \tag{3.8}
\]
Every source modulus is eventually below $N^{0.6}$.
Tested primes exceed $N/\log^{30}N$,
so $\Delta_0=\Delta_P$ on the original source.
Its squarefree NC and grouped allowed-repeat difference vanish before HB.
For $X_{Li}=\operatorname{Li}(Y_1)-\operatorname{Li}(Y_0)$, the term
$(m_A-X_{Li})/\varphi(q)$ is kept as the mass-conversion error.
The fixed logarithmic coefficient envelope and
$\sum_{q\le Q}\varphi(q)^{-1}\ll\log Q$ reduce it to PNT.
A sufficiently large fixed logarithmic saving gives $o(U_C)$.

Choose a tested-prime shell $[X,2X)$, with
$N/\log^{31}N\le X\le N$; on it HB has cap $2X^{1/3}$.
In each interior rectangle we choose the largest eligible elementary rough factor,
breaking ties by first index.
Its coefficient is rough M\"obius, rough unit, or normalized rough logarithm, of modulus at most one.
The complement has at most five factors and is bounded by $\tau_5(m)$.
With $\nu=N_{II}$, the full doubled support satisfies
\[
 X^{60\sigma}\le\nu\le X^{1/3+15\sigma},\quad M\nu\asymp X,
 \quad N\le X^{1+\sigma/20},\quad
 |mn-N|\ge\epsilon_A N/2
 \tag{3.9}
\]
This uses the original window $[N^{60\sigma},N^{1/3+14\sigma}]$ and
$\log_N(N/X)\le\sigma/1000$ eventually.
The shift stays $N$.

Lemma~\ref{lem:s2-preparation} supplies the equal-ratio mesh and integer endpoints.
For all four rows, including whole-$K$, it gives
$N_I\ge(2\nu)^{1-2^{-126}}$, $N_I\ge QX^{\sigma/3}$ and $M_I\ge X^{1/2}$.
It also gives H6, H7 and Class I.
The internal packing Lemma~\ref{thm:packing-input} applies at the literal $\Theta_b$,
and since $t_1<\beta<t_{\rm dd}$,
\eqref{p2:eq:row-level-floor} gives $\Theta_b\ge(1+t_1)/2$.
Put $T=t_1+2t_2\le3\beta$.
By (3.3), $t_3\le(K-T)/(4+2\tau)$, so
{\small
\[
\begin{aligned}
\ell_c=\Theta_b-\eta_L-t_1-t_2-t_3
&\ge\frac12-\eta_L-\frac T2-\frac{K-T}{4+2\tau}\\
&\ge\frac12-\eta_L-\frac{3\beta}{2}
              -\frac{K-3\beta}{4+2\tau}
>0.12953>\ell_0=\frac{53}{412}-\frac{\eta_L}{2}.
\end{aligned}
\]
}
Here $1/2-1/(4+2\tau)>0$ and $0<\tau=\sigma^9\le10^{-90}$.
Every selected S2-3 cell therefore has
$L_b^{\sigma^2}\ge N^{\ell_0\sigma^2}$.
Including slot-1 smooth inflation gives
\[
\begin{split}
 Q_1&<\nu X^{-\sigma},\\
 \nu^2Q_2Q_3^2&<X^{1-8\sigma},\\
 A_{\max}^{\theta}\nu Q_1Q_2^5Q_3^2&<X^{2-15\sigma},\\
 (\nu Q_3/Q_2)^{2\theta}\nu Q_1Q_2^5Q_3^2&<X^{2-15\sigma},
 \qquad \theta=7/32,\quad A_{\max}=X^{1+\sigma/20}.
\end{split}\tag{3.10}
\]
The ambient margins are
$(7271/256,52,95221/1280,22407/256)\sigma$.
After scale transfer, dyadic constants and endpoint localization, they exceed
$(28,51,74,87)\sigma$ eventually.
PF and the boundary region retain these margins and $\Theta_b$.

The dispersion argument is the one in
\hyperref[p2:dispersion-proof]{the distribution-input subsection}, with
\cite[Theorem~2.8, (3.17)--(3.20), (3.25)--(3.27), (3.57)--(3.60), (3.70)--(3.73)]{Z1}.
Its S2 parameters are as follows.
\begin{longtable}{@{}p{0.28\textwidth}p{0.67\textwidth}@{}}
\toprule
Step & Data \\
\midrule
First dispersion & $\sum_{m\in I_M}|\beta_m|^2\ll M_I\log^cX$.
For $\alpha^{(r)}=\alpha\mathbf1_{(n,r)=1}$, split at
$\|\alpha^{(r)}\|_2^2=\nu_I\log^{-A_3-1}X$. \\
SW precision & Below: trivial BDH.
Above: $\|\alpha^{(r)}\|_2(2\nu)^{1/2}\gg\nu\log^{-(334+A_3+1)/2}X$.
Request this extra saving before norm-normalized BFI. \\
Frequencies & $\eta_0=\sigma/1000$,
$H_2=QR^2S^2X^{\eta_0}/M_I$, $H_*=40H_2$;
$H_0<|h|\le\min(2H_0,H_*)$. \\
Mellin weight & In \eqref{p2:eq:dispersion-mellin}, set $a=N$.
Take $w=1$ on $[1/2,2]$, supported in $(1/4,4)$.
For $U(t)=m(t)$ and $L_{\rm osc}=(\log X)^{365}X^{\eta_0}$,
$U(t)\ll\min\{1,L_{\rm osc}^2/(1+|t|)^2\}$ and
$\int_{\mathbb R}U(t)\,dt\ll L_{\rm osc}\ll_\sigma X^{3\sigma/40}$.
Here $N/X\le(\log N)^{31}$. \\
Coloring & $d_{\rm disp}=1$, $\ell_0=53/412-1/20000$,
$u_0=N^{\ell_0\sigma^2}$, $\mathcal P_X=[u_0,X^3]\cap\mathbb P$.
Normalized cost: $2^{K_s+K_J}=O_\sigma(1)$. \\
Derivatives & $\epsilon_0=\sigma/2000>0$. \\
\bottomrule
\end{longtable}
Both norm cases give $\ll\nu_I\nu\log^{-A_3}X$.
The second dispersion retains
\[
 (s_1,Nc)=1,\qquad(s_1,q_{\rm disp}n_2)=1.
 \tag{3.11}
\]
Representations are normalized by $U(t)>0$, which is then restored.
At $U(t)=0$ the integrand vanishes.
Both nonmodular Mellin phases of \eqref{p2:eq:dispersion-mellin} are kept,
and coloring precedes aggregation.
Fiber multiplicity is at most $\tau(b)\tau_3(c)$. Hence
\[
 |A_\omega(t)|\ll_\sigma\frac{U(t)}{QR^2S^2},\qquad
 \sum_{b,c,f}|A^\chi_{b,c,f}(t)|^2
 \ll_\sigma U(t)^2X^{o(1)}\frac{\nu^2}{Q^2R^2S^3}.
 \tag{3.12}
\]
The spectral coefficient has the weighted norm
\[
 \sum_{n''\mid N^\infty}(n'')^\theta
 \|\beta^\flat(n'')\|_2^2
 \ll X^{o(1)}(H_2S)^2F,\qquad F\ll\nu/Q.
 \tag{3.13}
\]
For $y=f(h_1s'_1-h_2s_1)\ne0$, specialize to
\[
\begin{gathered}
 (r,s,d,c,n)=(z,p_0,b,c,p_0|y|),\quad
 a_T=\operatorname{sgn}(y)N,\quad(zb,p_0c)=1,\\
 (R_{sp},S_{sp},D_{sp},C_{sp},N_{sp})
 =(Z,\max(p_0,3)/2,B,C,8p_0Y),\\
 Z\in\{S^2,2S^2\},\quad B\ll\nu R,\quad
 C\ll\nu RS,\quad Y\ll H_2SF.
\end{gathered}
\]
with the least prime $p_0\nmid N$.
Then $p_0\ll\log N<z_0(X)$ and $p_0\nmid a_T$, and (3.13) gives the required weighted norm.
Here $p_0\mid c$ remains allowed.

For unit slot 2, its atom $s=1$ has $S=1/2$.
By (3.10), $H_2\ll X^{-9\sigma+\eta_0+o(1)}$,
so $H_*<1$ eventually and the frequency sum is empty.
Every other nonempty spectral scale is at least one.
For general $S$, $H_2/S\le X^{-9\sigma+\eta_0+o(1)}<1$,
hence $1\le H_*\le S$ eventually on nonempty frequency blocks.
Also $F\ge1/2$, as required at the weighted-norm endpoints.
Finally,
$A_{\max}C_{sp}D_{sp}N_{sp}R_{sp}S_{sp}
 \ll X^{\sigma/20+\eta_0+o(1)}\nu^4R^4S^6<X^{10.01}.$
By \cite[(3.89)]{Z1}, the positive-epsilon prefactor is eventually at most $X^{0.19019\sigma}$.
Equations (3.10)--(3.13) give the signed spectral bound.

Put $C_{\rm box}=C_b=334$ and $K_{\rm scale}=31$.
Lemma~\ref{lem:s2-preparation} gives the complete support
\[
 q<N^{e_D},\qquad e_D=153/256-\eta_L/2<0.6,
 \qquad q\le X^{0.6}\quad(N/\log^{31}N\le X\le N).
\]
This support includes repeated-prime moduli and every smooth divisor.
The full hypothesis check is as follows.
\begin{longtable}{@{}p{0.28\textwidth}p{0.67\textwidth}@{}}
\toprule
Condition & S2 data \\
\midrule
H1 & Coefficient (3.6), with primes avoiding $2N$; corrections exactly (3.8). \\
H2 & $d_{AA}$ lies in slot 1; slots 2 and 3 are $u_0$-rough.
Keep repeated primes and empty slots. \\
H3 and LEN & (3.9)--(3.13), the full support above and
Lemma~\ref{lem:s2-preparation}.
Majorant: $m_0=M(1+r_X^r)/2$, $1\le r\le5$,
$|m_0|\le2M\ll X/\nu$.
Ambient length: $N_I\ge(2\nu)^{1-2^{-126}}$. \\
H4 & $a=N$, $(q,N)=1$;
$0<N\le X\log^{31}N\le X^{1+\sigma/20}$. \\
H5 & $N^{60\sigma}\le\nu\le N^{1/3+14\sigma}$ and (3.9),
with Lemma~\ref{lem:s2-preparation}'s endpoints. \\
H6 & $E_{\log}\ll_\sigma N(\log N)^{18-C_b}$;
$E_{\rm bdry}\ll_\sigma N(\log N)^{16-C_b/2}$. \\
H7 & (3.16), using external coefficients only. \\
\bottomrule
\end{longtable}
Theorem~\ref{thm:distribution-input}, or
\cite[Theorem~3.2, (3.13)]{Z1}, therefore applies.
The complete-weight prime-source theorem has a different coefficient family.
\end{proof}

\subsection{Shiu's estimate and the three dispersion terms}
\label{sec:s2-4}

\begin{lemma}
\label{lem:external-classii-sum}
Let $\mathscr F$ be a prepared family of interior Class-II boxes satisfying
the local scale, roughness, support and coprimality hypotheses of
Theorem~\ref{thm:distribution-input} and (3.9)--(3.11), with uniform fixed
source constants. Assume the coefficient and scale bounds (3.16), with
$\mathcal J$ replaced by $\mathscr F$ and $K_0=8042$. Then (3.15) holds
with a fixed exponent $C_{\mathscr F}$ in place of $C_{S2}$, and the
normalized sum satisfies (3.17).
\end{lemma}
\begin{proof}
We first bound the deletion of $f_{\rm gcd}=(m,r_2)>z_0$.
The restoration formula in \cite[\S3.5, (3.45)]{Z1} retains $(m,r_2)=1$
by the exact M\"obius expansion.

Let $MN\asymp x$ and $L=\log x$, with local shift $a$.
Assume fixed $0<\delta<1$, $0<\varepsilon<1/2$, and
\[
 x^{2\delta}\le N\le x^{1/3+\delta/2},\quad
 |I_N|=N_I\asymp NL^{-C_b},\quad
 |\widetilde I_M|=2M_I\asymp ML^{-C_b},
\]
with $I_N\subset[N,2N]$ and $\widetilde I_M\subset[M/4,3M]$.
On the doubled box assume $|mn-a|\ge xL^{-C_b}\ge1$, and finally assume $0<|a|\le x^{1+\varepsilon}$ and $1\le f_{\rm gcd}\le Mx^{-\delta}$.
For fixed nonnegative $B_i$,
\[
 \sum_{n\in I_N}\tau(n)^{B_1}
 \sum_{\substack{m\in\widetilde I_M\\f_{\rm gcd}\mid m}}
 \tau(m)^{B_2}\tau(|mn-a|)^{B_3}
 \ll\frac{M_IN_I}{f_{\rm gcd}}(\tau(f_{\rm gcd})L)^{O_{B_i}(1)}.
 \tag{3.14}
\]
To prove this, write $m=f_{\rm gcd}m'$, so the interval length is $H'=2M_I/f_{\rm gcd}\gg x^{\delta/2}$.
For fixed $n$, put $g=(f_{\rm gcd}n,|a|)$,
$v=f_{\rm gcd}n/g$ and $a'=a/g$.
Then $(v,a')=1$, $|mn-a|=g|vm'-a'|$.
Put $\ell(m')=f_{\rm gcd}nm'-a$; consecutive values differ by
$f_{\rm gcd}n\le(Mx^{-\delta})(2N)\ll x^{1-\delta}$.
Opposite signs would imply
$|\ell(m')|+|\ell(m'+1)|=f_{\rm gcd}n\ge2xL^{-C_b}$, a contradiction, so the sign stays constant.
The values $|vm'-a'|$ lie in one reduced class modulo $v$.
Their interval has length $Y=vH'$ and endpoint $Z\ll x^{1+\varepsilon}$.
Choose fixed $0<b<\min(1/4,\delta/[8(1+\varepsilon)],\delta/8)$.
Since $v\ll x^{1+o(1)}$ and $H'\gg x^{\delta/2}$,
$Y\ge Z^b$ and $v\le Y^{1-b}$ eventually.
Shiu's reduced-progression bound \cite{Shiu1980} gives
$\sum_{m'}\tau(|vm'-a'|)^C
 \ll\frac{Y}{\varphi(v)}L^{2^C-1}\ll H'L^{O_C(1)}.$
For modulus one,
$\sum_{m'}\tau(f_{\rm gcd}m')^C\ll\tau(f_{\rm gcd})^CH'L^{O_C(1)}.$
We use Cauchy--Schwarz and $\tau(g)\le\tau(f_{\rm gcd})\tau(n)$, then the localized divisor moment.
This proves (3.14), including signed shifts.
The modulus condition uses $Y=vH'$ and $Y/v=H'$.

Take $x=X$, $N_{\rm local}=\nu$, $\delta=30\sigma$,
$\varepsilon=\sigma/20$ and $a=N_{\rm ambient}$.
By \cite[(3.6)]{Z1} and (3.9)--(3.10),
$f_{\rm gcd}=(m,r_2)\le r_2\ll M_I X^{-89\sigma}\ll M X^{-30\sigma},$
where $M_I\ll M(\log X)^{-334}$ and $\sigma>0$ is fixed.
Thus (3.14) applies.

In the following divisor sum, $r_2$ is the second two-rough-slot product
of \cite[(3.28)--(3.31)]{Z1}.
Its primes are at least $u_0=N^{\ell_0\sigma^2}$ and $r_2\le X^{0.6}$.
Counting repeats gives
$\Omega(r_2)\le\frac{0.6}{\ell_0\sigma^2},\qquad
 \tau(r_2)\le2^{\Omega(r_2)}\ll_\sigma 1.$
This includes $r_2=1$.
Every nontrivial $f_{\rm gcd}\mid r_2$ satisfies
$f_{\rm gcd}\ge u_0>z_0(X)$ eventually.
For fixed $B\ge0$,
$\sum_{\substack{f_{\rm gcd}\mid r_2\\f_{\rm gcd}>z_0(X)}}\frac{\tau(f_{\rm gcd})^B}{f_{\rm gcd}}
 \le\frac{\tau(r_2)^{B+1}}{u_0}\ll_{\sigma,B}\frac1{u_0}.$
We partition by the exact value $f_{\rm gcd}=(m,r_2)$.
Dropping the remaining gcd condition and summing (3.14) gives
\[
 \sum_{n\in I_\nu}\tau(n)^{B_1}
 \sum_{\substack{m\in\widetilde I_M\\(m,r_2)>1}}
 \tau(m)^{B_2}\tau(|mn-a|)^{B_3}
 \ll_\sigma\frac{M_I\nu_I}{u_0}(\log X)^C,\qquad |I_\nu|=\nu_I.
\]
The divisor weights fix $C$ independently of $r_2$.

The mixed term keeps $qr_1\mid mn_1-a$.
For fixed $(m,n_1)$, the number of admissible pairs $(q,r_1)$ is at most
$\tau(|mn_1-a|)^2$.
Also $\varphi(qr_2)^{-1}\ll(\log X)^C/(Qr_2)$,
$|\lambda_{r_i}|\ll_\sigma1$ and $\sum_{n_2}1\le\nu_I$.
Increasing $B_3$ by two and using $\sum_{r_2\asymp RS}r_2^{-1}\ll1$, we find the positive deletion error
$E_{\rm gcd}\ll_\sigma\frac{M_I\nu_I^2}{Q u_0}(\log X)^C.$
For pure density,
$\sum_{q\asymp Q}\varphi(q)^{-2}\sum_{r_1\asymp RS}\varphi(r_1)^{-1}
\ll(\log X)^C/Q$.
This gives the same bound.
Copy symmetry gives the first-copy deletion bound.
For two congruences, $qr_2\mid mn_2-a$ and $(qr_2,a)=1$
imply $(m,r_2)=1$, so its deleted contribution is empty.
These three bounds include unit rough slots, and both masks (3.11) remain in force.

The first Cauchy--Schwarz inequality is
$|S|^2\ll QM_I(\log X)^C S'$.
Its full divisor-deletion contribution is at most
\[
 (QM_I(\log X)^CE_{\rm gcd})^{1/2}
 \ll_\sigma M_I\nu_I(\log X)^C u_0^{-1/2}
 \ll_\sigma X^{1-\ell_0\sigma^2/2}(\log X)^C.
\]
Since $\ell_0\sigma^2/2=2c_\sigma>c_\sigma$, this term is absorbed by
the power-saving term established below in (3.15).
The exact mixed term retains $(m,r_2)=1$ via \cite[\S3.5, (3.45)]{Z1},
and the divisor sum above gives its full alternative bound.
Restoration and near-diagonal/BDH recombination follow
\cite[(3.34)--(3.47), (3.49)--(3.53)]{Z1}, with both masks and all three dispersion terms kept.

Theorem~\ref{thm:distribution-input} is applied at $A_{\rm work}=A'+4$.
Its local terms are $X_j\log^{-A_{\rm work}}X_j$ and
$X_j^{1-c_\sigma}\log^{C_{Z1}}X_j$, where
$c_\sigma=\frac14(\frac{53}{412}-\frac1{20000})\sigma^2>0$.
The latter includes GR restoration and first-CS deletion.
Here $\ell_0=53/412-1/20000$, $\ell_0\sigma^2/2=2c_\sigma$.
Three internal modulus-slot shells cost at most $\log^3X_j$, so the first term is at most
$X_j\log^{-A'-1}X_j\le X_j\log^{-A'}X_j$.
With shell and source-specific GR logarithms absorbed into $C_{S2}$,
\[
 |R_j|\ll_\sigma X_j\log^{-A'}X_j
       +X_j^{1-c_\sigma}\log^{C_{S2}}X_j.
 \tag{3.15}
\]

For the external family $\mathcal J$, Lemma~\ref{lem:s2-preparation} gives
\[
 \sum_{j\in\mathcal J}|c_j|\le C_\sigma\log^{K_0}N,
 \qquad N/\log^{31}N\le X_j\le N,
 \tag{3.16}
\]
Here $K_0=K_0^{S2}$ is fixed before $A'$.
Put $L=\log N$.
Lemma~\ref{lem:direct-preparation} gives
$K_*=6(C_b+1)+2=2012$.
We use the larger envelope $K_0=24C_b+26$:
$K_0^{S2}=24\cdot334+26=8042,\qquad
 A'=8047,\qquad A_{\rm work}=8051.$
PF and the selected boundary region form one whole-$K$ family.
Their prime realizations and smooth divisors remain internal.

For $\mathscr F$, we replace $\mathcal J$ and $C_{S2}$ in (3.15)--(3.16)
by $\mathscr F$ and any fixed $C_{\mathscr F}$.
With $K_0=8042$, $A'=8047$ and denominator (3.1), we get
\[
 \begin{aligned}
 E_{II}(\mathscr F)&:=\frac1{W_N}\sum_{j\in\mathscr F}|c_j||R_j|\\
 &\ll_\sigma L^{K_0+2-A'}\\
 &\qquad+N^{-c_\sigma}L^{K_0+C_{\mathscr F}+2+31c_\sigma}\\
 &\ll_\sigma L^{-3}+N^{-c_\sigma}L^{8044+C_{\mathscr F}+31c_\sigma}=o(1).
 \end{aligned}
 \tag{3.17}
\]
Here $8044=K_0+2$; the extra 2 comes from $W_N\gg N/L^2$.
By (3.1), $W_N\gg N/L^2$ and
$\tfrac12\log N\le\log X_j\le\log N$.
The bound $X_j\ge N/L^{31}$ gives $31c_\sigma$.
Also $K_0+2-A'=8042+2-8047=-3$,
and for fixed $B$, $e^{-c_\sigma t}t^B\to0$ as $t=\log N\to\infty$.
Here $\sigma>0$ and all source and box parameters are fixed before $N_0$.
For S2 we take $\mathscr F=\mathcal J$, $C_{\mathscr F}=C_{S2}$ and
$E_{II,S2}=E_{II}(\mathcal J)$.
\end{proof}

\subsection{The strict BV complement}
\label{sec:s2-5}

\begin{lemma}
The source on $\mathcal B$ has upper main comparator (3.19)
and remainder $o(U_C)$, at the actual cutoff $p_3$.
\end{lemma}
\begin{proof}
For $r\in\mathcal B$, put
$b=p_1p_2p_3,\quad Q_{BV}=N^{1/2-\kappa}\log^{-B}N,
 \quad\kappa=10^{-6},\quad D_b=Q_{BV}/b,\quad z_b=\sqrt{D_b}.$
On all of $R1$,
\[
 \log_ND_b\ge3/206-\kappa-o(1)>0,\qquad
 2r_3-\log_ND_b\ge63823/309000+\kappa+o(1)>0.
 \tag{3.18}
\]
Write $L_N=\log N$. Since $\mathcal B\subseteq R1$,
$g=r_1+2r_2+4r_3>H=249/250$. Ordering gives $g\le7r_1$, so
$r_1>H/7$ and
\[
 r_1+r_2+2r_3=\frac{g+r_1}{2}>\frac{4H}{7}.
\]
The definition of $D_b$ now gives
\[
\begin{aligned}
 r_3-\log_ND_b
 &=r_1+r_2+2r_3-\frac12+\kappa
       +B\frac{\log\log N}{\log N}\\
 &>\frac{4H}{7}-\frac12+\kappa
       +B\frac{\log\log N}{\log N}\\
 &=\frac{121}{1750}+\kappa
       +B\frac{\log\log N}{\log N}>0.
\end{aligned}
\]
Hence $1<D_b<p_3$ for all sufficiently large $N$, and
$z_b=\sqrt{D_b}<p_3$. Sifting is monotone in its cutoff, so
\[
 S_{\rm pl}(b,p_3)\le S_{\rm pl}(b,z_b).
\]
Apply the $s=2$ upper sieve to the right-hand side.
Its squarefree weight has $|\lambda^+_{b,d}|\le1$,
$d<D_b$ and $d\mid P_N(z_b)$. Thus every prime factor of $d$
is below $p_3$; indeed $d<D_b<p_3$.
For $k=bd$, the three largest distinct prime factors recover $b$ uniquely, also when $d=1$.
The actual aggregate weight
$a_k=\lambda^+_{b,d}\mathbf1_{\{b\text{ belongs to }\mathcal B\}}$
has $\|a\|_\infty\le1$, $\|a\|_2^2\le Q_{BV}$, and the prefix/divisor multiplicity is one.
Here $N\bmod k$ is reduced.
The two-endpoint remainder is at most twice the classical BV maximal sum.
It is $O_A(N\log^{-A}N)=o(U_C)$ for sufficiently large fixed $B$.
The positive upper comparator is enlarged from $\mathcal B$ to $R1$:
\[
 I_{BV}=\int_{R1}\frac{4\,dr_1dr_2dr_3}
 {r_1r_2r_3(1/2-\kappa-r_1-r_2-r_3)}.
 \tag{3.19}
\]
The factor 4 includes Mertens and $F(2)$.
Thus $I_{BV}$ enters directly, with its stated normalization.
The comparator overlaps $\mathcal F\cap R1$.
The source partition (3.4) stays disjoint, with its original boundary conventions.
\end{proof}

\subsection{The signed functional, tail and endpoint}
\label{sec:s2-6}\label{sec:s2-7}

\begin{lemma}
The selected endpoint is (3.22).
With its tail allowance included, the source satisfies (3.23).
\end{lemma}
\begin{proof}
Put $a=2/e^\gamma$. Define
\[
 v(t)=\begin{cases}(1+t)/2,&0\le t\le25/128,\\
 57/64-3t/2,&25/128<t\le1/3,\\0,&\text{otherwise},\end{cases}
 \quad\psi(x,y)=x\mathbf1_{\{x\ge y\}}.
\]
Let $\vartheta(r_1,r_2,r_3)$ be the maximum of
\[
\begin{gathered}
 v(r_1),\ v(r_2),\ v(r_1+r_2),\ v(r_1+r_2+r_3),\\
 \psi(\min(153/256-3r_3/8,1-2r_2),(1+r_1)/2),\\
 \psi(\min(153/256-3r_2/8,1-2r_3),(1+r_1)/2),\\
 \psi(v(r_1+r_3),r_1+2r_2+r_3),\quad
 \psi(v(r_2+r_3),2r_1+r_2+r_3).
\end{gathered}
\]
Put
\[
\begin{split}
 s_0&=19037/64,\quad
 a_1(r)=(\vartheta(r,u,u)-u-r)/u,\\
 a_2(r_1,r_2)&=(\vartheta(r_1,r_2,u)-u-r_1-r_2)/u,\\
 s_3(r)&=(\vartheta(r)-u-r_1-r_2-r_3)/r_3,\\
 P&=\frac au\left[f(s_0)-\int_u^\beta F(a_1(r))\frac{dr}{r}
 +\int_{u\le r_2<r_1<\beta}f(a_2(r_1,r_2))\frac{dr_1dr_2}{r_1r_2}\right],\\
 d\mu&=\frac{dr_1dr_2dr_3}{r_1r_2r_3^2},\quad
 V=\int_\Delta d\mu,\quad V_B=\int_{R1}d\mu,\\
 J_{PF}&=\int_{PF}[F(\max(s_3(r),1))-1]d\mu.
\end{split}
\]
The signed functional is
\[
 T8=P-aV-aJ_{PF}+(aV_B-I_{BV})
   =P-a\int_{PF}F(\max(s_3,1))d\mu-I_{BV}.
 \tag{3.20}
\]
It combines the prefix terms with the third-term comparison.
On $PF$, $\vartheta(r)\ge(1+r_1)/2$ and $g(r)\le H=1-2u$. Thus
$\vartheta(r)-u-r_1-r_2-2r_3\ge\frac{H-g(r)}2\ge0.$
Hence $s_3(r)\ge1$ and $F(\max(s_3,1))=F(s_3)$ on PF, with equality possible on its boundary.
The analytic whole-$K$ argument $1.0117$ uses level loss $\eta_L$.
At depth one and $u\le r<\beta$, the $w$ guard is active and its
branch gives $\vartheta(r,u,u)\ge19101/32000$ throughout the whole
parent cell. Thus the numerator margin is at least
$u-\eta_L=19/10000>0$.

The source/profile numerator margins are
$u-\eta_L-9\beta\tau>0$ on triples and
$u/2-\eta_L-5\beta\tau>0$ on second prefixes.
By continuity and the 2-Lipschitz S2 hull bound, and since $v(r_1)$ dominates at indicator switch-off,
$\Theta_b\ge\vartheta(r)-6\beta\tau$.
Inherited endpoints cost at most another $3\beta\tau$,
and the numerical and analytic losses are $u$ and $\eta_L$.
Missing-coordinate infima retain the floor $153/256-3u/8$
and the relevant $v(r_1+r_2)$ candidate.
Depth-two padding costs at most $u/2+5\beta\tau$.
Increasing $f$ and decreasing $F$ give the signs in (3.20).
PNT/Mertens, harmonic prime measures and equal-bin factorials give the ordered integrals.
Here $\tau$ is fixed throughout.

The boundary region has width
$\delta_K=\rho K-H=\tau(H+2\beta)+2\beta\tau^2.$
Its outer ordered area is at most $\beta^2/2$, and each fiber has width at most $\delta_K/4$.
Also $aF(1)=4$.
Including the dynamic cutoff,
\[
 E_{spill,D}\le\frac{(1+2\tau)\beta^2\delta_K}{8u^4}
 <275000000\tau,\qquad E_{dyn,D}\le2\tau V<9680\tau.
 \tag{3.21}
\]
Thus the cell-selection/dynamic allowance is $275009680\tau$,
while the boundary arithmetic discrepancy is estimated separately.

The directed 192-bit Arb inputs \cite{Johansson2017} are:
\begin{center}
\begin{tabular}{@{}ll@{}}
Prefix & $n=600$ \\
PF upper & $n=1536$ \\
Corner lower & $n=2800$
\end{tabular}
\end{center}
The true sieve functions $f,F$ are used beyond 12.
Put $E(s)=\max(F(s)-1,1-f(s))$.
The finite delay-integral enclosure gives $E(s)<6\cdot10^{-6}$ on $[11,12]$.
The classical linear-sieve asymptotic gives
$F(s)-1=O(e^{-s})$ and $1-f(s)=O(e^{-s})$.
Thus $e^sE(s)$ is bounded for all sufficiently large $s$.
Since $f,F$ are continuous, it is bounded on the remaining compact interval.
Hence $B=\sup_{s\ge11}e^sE(s)<\infty$ before the contraction below.
The same decay gives integrable tails and the vanishing limits used next.
Integrating $(s(F(s)-1))'=f(s-1)-1$ and
$(s(1-f(s)))'=1-F(s-1)$, with their vanishing limits at infinity,
gives the exact identities
\[
 s(F(s)-1)=\int_{s-1}^\infty(1-f(t))dt,\qquad
 s(1-f(s))=\int_{s-1}^\infty(F(t)-1)dt
\]
and hence, for $s\ge12$,
\[
 sE(s)\le\int_{s-1}^\infty E(t)dt\le Be^{1-s}.
\]
Separating $11\le s\le12$ from $s\ge12$ gives
$B\le\max(6\cdot10^{-6}e^{12},eB/12)$.
Since $e/12<1$, this implies
$B\le6\cdot10^{-6}e^{12}<0.976529<1$.
The preceding integral bound therefore gives
$E(s)\le e^{1-s}/s$ for $s\ge12$.

Fix $u\le r_2\le r_1\le\beta$ and write $R=r_1+r_2$, $x=r_3$,
$C=153/256$, $L=57/64$ and $A=(1+r_1)/2$.
The PF fiber has $u\le x<r_2$ and $r_1+2r_2+4x\le H$.
On every affine branch $\phi(x)$ in the hull, its argument is
$s_\phi(x)=c/x+k$, where
$s_\phi=(\phi(x)-u-R-x)/x$.
The following list covers every branch of $v(r_1)$, $v(r_2)$,
$v(R)$, $v(R+x)$, the two unguarded parts of each of the first two
$\psi$ terms, and both branches of each pair guard:
\begingroup\scriptsize
\setlength{\tabcolsep}{1.5pt}
\begin{longtable}{@{}p{0.22\textwidth}p{0.43\textwidth}p{0.26\textwidth}c@{}}
\toprule
Branch $\phi$ & $c$ & Lower bound for $c$ & $k$ \\
\midrule
$v(r_1)$ & $1/2-u-r_1/2-r_2$ & $1/2-u-3\beta/2$ & $-1$ \\
$v(r_2)$ & $1/2-u-r_1-r_2/2$ & $1/2-u-3\beta/2$ & $-1$ \\
$v(R)$, increasing & $1/2-u-R/2$ & $1/2-u-\beta$ & $-1$ \\
$v(R)$, decreasing & $L-u-5R/2$ & $L-u-5\beta$ & $-1$ \\
$v(R+x)$, increasing & $1/2-u-R/2$ & $1/2-u-\beta$ & $-1/2$ \\
$v(R+x)$, decreasing & $L-u-5R/2$ & $L-u-5\beta$ & $-5/2$ \\
$q_1=C-3x/8$ & $C-u-R$ & $C-u-2\beta$ & $-11/8$ \\
$q_1=1-2r_2$ & $1-u-r_1-3r_2$ & $1-u-4\beta$ & $-1$ \\
$q_2=C-3r_2/8$ & $C-u-r_1-11r_2/8$ & $C-u-19\beta/8$ & $-1$ \\
$q_2=1-2x$ & $1-u-R$ & $1-u-2\beta$ & $-3$ \\
$v(r_1+x)$, increasing & $1/2-u-r_1/2-r_2$ & $1/2-u-3\beta/2$ & $-1/2$ \\
$v(r_1+x)$, decreasing & $L-u-5r_1/2-r_2$ & $L-u-7\beta/2$ & $-5/2$ \\
$v(r_2+x)$, increasing & $1/2-u-r_1-r_2/2$ & $1/2-u-3\beta/2$ & $-1/2$ \\
$v(r_2+x)$, decreasing & $L-u-r_1-5r_2/2$ & $L-u-7\beta/2$ & $-5/2$ \\
\bottomrule
\end{longtable}
\endgroup
Here $q_1=\min(C-3x/8,1-2r_2)$ and
$q_2=\min(C-3r_2/8,1-2x)$.
Substitution of $u=1/500$ and $\beta=50/309$ shows that every lower
bound in the table is at least
$c_{\min}=L-u-5\beta=196681/2472000>0$.
The permanent baseline $A=v(r_1)$ is always present.
The guards $q_i\ge A$ switch off only, since each $q_i$ is continuous
and nonincreasing in $x$; at switch-off its value equals $A$.
For the pair guards put
\[
 D_1=v(r_1+x)-(r_1+2r_2+x),\qquad
 D_2=v(r_2+x)-(2r_1+r_2+x).
\]
Since $r_i+x\le2\beta<1/3$, these functions are continuous,
strictly decreasing, with slopes $-1/2$ and $-5/2$ on the two branches.
On the increasing branch $D_i\ge1/2-3\beta=3/206>0$,
so switch-off can occur only on the decreasing branch.
At a switch-off, direct substitution gives
\[
\begin{aligned}
 (r_1+2r_2+x)-A&=-r_1/2+6r_2/5-23/160
                  \le-1507/49440<0,\\
 (2r_1+r_2+x)-A&=7r_1/10-23/160
                  \le-1507/49440<0.
\end{aligned}
\]
Thus removing either guarded branch cannot raise the hull.
All affine crossings and the kink of $v$ are continuous.
The only other drop is $v(R+x)$ at $R+x=1/3$;
its left value is $25/64<A$, so it is inactive there.
Every active branch has $s_\phi'(x)=-c/x^2\le-c_{\min}/x^2$.
The full argument hull $s_3$ is therefore continuous and strictly decreasing on the fiber.
Its finitely many affine pieces have disjoint $s$-images.
Consequently, with no branch-count factor,
\[
\begin{aligned}
 \int_{\{x:s_3(x)>12\}}(F(s_3(x))-1)\frac{dx}{x^2}
 &\le\frac1{c_{\min}}\int_{12}^{\infty}E(s)\,ds\\
 &\le\frac{e^{-11}}{12c_{\min}}.
\end{aligned}
\]
Integrating over the full ordered $(r_1,r_2)$ triangle gives
\[
 \int_{u\le r_2<r_1\le\beta}\frac{dr_1dr_2}{r_1r_2}
 =\frac12\log^2(\beta/u).
\]
The true PF tail above 12 costs less than
\[
 \frac a4\frac{\log^2(\beta/u)}2\frac{e^{-11}}{12c_{min}}
 <4.739184885736468\cdot10^{-5}\quad\text{D units}.
\]
Here $a=2e^{-\gamma}>0$.
The lower endpoint uses $P_{lo}$, upper products $aV$ and
$a(J_{PF})_{hi}$, and the lower corner endpoint.
All prefix arguments exceed 128.
Their repair costs less than $1.154\cdot10^{-54}$ D.
Both repairs fit strictly inside $1/20000$.
Adding the true tail to the finite PF upper bound gives
\[
 \begin{aligned}
 d_{T8}&=\frac{P_{lo}-aV-a(J_{PF})_{hi}+(aV_B-I_{BV})_{lo}}4
                 -\frac1{20000}\\
 &>1.61958135698908352559133404332105729220076714081650.
 \end{aligned}
 \tag{3.22}
\]
This eight-branch bound is in D units.
It includes that allowance.

The source identities and preceding estimates yield
\[
 \frac{S_{2,\rm raw}}{W_N}
 \ge d_{T8}-275009680\tau-E_{AA,S2}-E_{II,S2}-e_{rest,S2}(N),
 \qquad e_{rest,S2}(N)\longrightarrow0.
 \tag{3.23}
\]
The four rows and three third-row pieces form \texttt{S2-T8/4}.
Equation (3.17) includes CRT, colors, zero mode, Poisson, Theorem 2.8
and all three GR/dispersion terms.
At $b=1$, the source-square restoration is empty.
The weighted low-tested-prime and exceptional-source bounds are
$O_u(\log^{-28}N)$ and $O_u(\log^3N/N)$.
Direct proper prime powers have raw bound $X^{1/2+o(1)}$, and their actual scalar sums are retained.
Theorem~\ref{thm:repeated-primes} bounds the forbidden square union by
$O_\sigma(N^{-\ell_0\sigma^2}\log^4N)$.
The NC and allowed-repeat terms vanish in the actual-mass convention of (3.8).
The Li convention retains the weighted source mass conversion.

For H6, the normalized LOG and boundary errors are
$O_\sigma(L^{20-C_b})$ and $O_\sigma(L^{18-C_b/2})$.
Class I is $O_\sigma(L^{-4})$, with $A_I=2018$, and at $C_b=334$ these enter $e_{rest,S2}$.
The strict BV source is $\mathcal B$, with main term in (3.20).
PNT, Mertens, endpoint and mass-conversion errors vanish.
The fixed $275009680\tau$ lies outside $d_{T8}$, while the allowance $1/20000$ lies inside it, in D units.
Finally, \texttt{E\_AA,S2} is the S2 share of the global AA37/G2 error.
Its fixed $O(\sigma)$ part stays separate from $e_{rest,S2}$.
Each contribution enters once.
The joint functional includes the S2 $f$-to-$F$ comparison and R1 corner.
The S1 $f$-to-$F$ and G1 PR-M terms remain separate.
This proves (3.23), which Section~6.3.5 uses in the full error sum.
\end{proof}
